\section{Introduction}

Many complex stochastic systems exhibit fluctuations that are neither
independent nor characterized by a single scale. Financial returns, turbulent
cascade surrogates, and other intermittent signals often display heavy-tailed
increments, clustered activity, long-range dependence in amplitudes, and
nonlinear scaling of empirical moments. Multifractal analysis provides a
compact language for such behavior by describing how moments of different
orders scale across observation scales. In this language, a process is not
summarized only by a single Hurst or roughness exponent, but by a scaling
function \(\zeta(q)\) and, in the broader formalism, by an associated
singularity spectrum \(f(\alpha)\)
\cite{halsey1986fractal,chhabra1989direct}. This viewpoint has been widely
used in complex systems and econophysics, where apparent multiscaling is often
interpreted as a signature of cascade-like organization or intermittent
volatility dynamics
\cite{cont2001stylized,mandelbrot1997multifractal,calvet2002multifractality,dimatteo2005long}.

Among the models used to represent such spectra, the multifractal random walk
(MRW) is particularly attractive because it gives a low-dimensional analytic
description of nonlinear scaling. In the convention used in this work, the
MRW-like spectrum is written as
\[
\zeta_{\mathrm{MRW}}(q)=qH-\frac{1}{2}\lambda^2 q(q-2),
\]
where \(H\) controls the leading monofractal scaling and \(\lambda^2\) controls
the curvature or intermittency of the spectrum
\cite{bacry2001multifractal,bacry2003log}. When \(\lambda^2\) is close to zero,
the spectrum is nearly linear and resembles a monofractal scaling law such as
that of fractional Brownian motion or fractional Gaussian noise
\cite{mandelbrot1968fractional}. When \(\lambda^2\) is larger, the spectrum
becomes visibly curved and is more compatible with an MRW-like interpretation.
This simple analytic form makes MRW a useful reference model for finite-sample
spectral diagnostics.

The central difficulty is that empirical multifractal spectra are not observed
directly. They are estimated from finite time series through structure
functions, detrended fluctuation statistics, wavelet-transform modulus maxima,
wavelet leaders, or related multiscale estimators
\cite{muzy1991wavelets,muzy1993multifractal,arneodo1995thermodynamics,kantelhardt2002mfdfa,jaffard2007wavelet,wendt2007bootstrap}.
These
estimators are sensitive to the available sample length, the scale range, the
chosen q-grid, and the stability of high-order moments. A finite monofractal
sample can generate apparent curvature; heavy-tailed data can destabilize
high-\(q\) moments; and regime-switching or volatility-clustering signals can
produce scaling patterns that resemble multifractal behavior without implying
an MRW mechanism
\cite{bouchaud2000apparent,ludescher2011finite,grech2013spurious,zhou2009components,barunik2012source}.
Thus the practical problem is not merely to fit a curved spectrum, but to
decide whether the curvature is stable, whether a monofractal explanation is
already sufficient, and whether the finite sample contains enough information
to support an MRW-like interpretation.

Traditional multifractal estimators are indispensable, but they typically leave
this decision to post-hoc interpretation. A structure-function regression may
produce an empirical \(\zeta(q)\), and MFDFA or wavelet-based methods may
produce more robust multiscale estimates, but the analyst must still decide
whether the observed curvature is physical, finite-size induced, tail-driven,
or a consequence of scale selection. Even for wavelet-based multifractal
procedures, reliable curvature estimation depends on the asymptotic regime and
the available scale support \cite{mallat1992singularity,bacry2010mixed}.
Conversely, unconstrained neural networks
can be trained to regress parameters or spectra from finite signals, but a
black-box prediction does not by itself enforce spectral geometry, distinguish
monofractal and MRW explanations, or expose when the inverse problem is
under-identified. This is precisely the setting where constrained scientific
machine learning is useful: data-driven components should be combined with
mathematical structure, physically interpretable coordinates, and explicit
diagnostics rather than used as unrestricted replacements for the model
\cite{arridge2019inverse,cranmer2020sbi,raissi2019pinn,karniadakis2021piml}.
This places the neural component of CMIN-SR closer to data-driven inverse
modeling and simulation-based diagnostic inference than to unrestricted
time-series prediction. Related scientific machine-learning work makes a
similar distinction between using neural networks as unconstrained predictors
and using them as structure-aware components in physical modeling pipelines
\cite{brunton2020mlfluid,karniadakis2021piml}. The term
simulation-based is used here in this broad methodological sense; the present
paper does not perform posterior inference or claim a full likelihood-free
Bayesian analysis.

This paper introduces CMIN-SR, a validity-aware spectral diagnostic framework
for finite-sample MRW inference. The aim is not to claim guaranteed recovery of
the true MRW intermittency parameter from short observations. Instead, the
framework separates the problem into three interpretable stages. First, a
finite signal is mapped to an empirical scaling spectrum
\(\zeta_{\mathrm{emp}}(q)\). Second, this empirical spectrum is projected onto
two competing analytic families: a linear monofractal spectrum and a parabolic
MRW spectrum. Third, the resulting residuals, curvature diagnostics, boundary
scores, and instability indicators are passed to a spectral-geometry
calibrator that reports diagnostic scores for scaling stability, nonlinear
curvature, monofractal compatibility, MRW compatibility, and low-curvature
boundary behavior. In this formulation, \(\lambda^2_{\mathrm{proj}}\) is a
projection coordinate within the MRW family, not proof that the observed signal
was generated by an MRW process.

This diagnostic formulation leads to a different scientific question from
ordinary parameter regression. Rather than asking only whether a neural model
can output a number close to \(\lambda^2\), we ask whether the finite empirical
spectrum contains stable evidence that favors a curved MRW explanation over a
linear monofractal one. This distinction is important because stable scaling
alone is not sufficient evidence for multifractality. It is also important
because failure can occur at different levels: the analytic spectral geometry
may be interpretable, while the raw empirical spectrum may be too noisy or too
short to preserve the relevant curvature.

The contributions of this work are as follows.

\begin{enumerate}
  \item We formulate finite-sample MRW inference as a validity-aware spectral
  diagnostic problem. The framework explicitly separates empirical spectrum
  estimation, monofractal/MRW projection, residual geometry, curvature
  diagnostics, boundary behavior, and instability warnings.

  \item We introduce and evaluate a spectrum-space calibration layer that
  operates on analytic and semi-analytic \(\zeta(q)\) curves. In clean spectrum
  space, the calibrator separates linear monofractal spectra, low-curvature
  boundary MRW spectra, strongly curved MRW spectra, and unstable distorted
  spectra, demonstrating that the spectral interpretation layer is learnable
  when the input geometry is controlled.

  \item We provide a finite-sample identifiability analysis showing that the
  main bottleneck lies upstream, in empirical spectrum estimation. Under the
  tested short-window settings, both deterministic structure-function
  estimators and neural empirical-spectrum estimators struggle to recover MRW
  curvature reliably. This negative result is central to the paper: it
  motivates the interpretation of CMIN-SR as a calibrated diagnostic framework
  rather than a guaranteed short-window MRW parameter recovery engine.
\end{enumerate}

The remainder of the paper is organized as follows. Section~2 introduces the
finite-sample spectral inference problem, the monofractal and MRW projection
families, and the diagnostic inverse formulation. Section~3 describes the
CMIN-SR framework and shows how MRW samples are simulated. Section~4 presents
the experimental design, including analytic spectra, synthetic process
families, classical baselines, zeta-space noise, failure attribution, and a
cautious real-data sanity check. Section~5 reports the main results.
Section~6 analyzes failure modes, and Section~7 discusses interpretation and
limitations.
